\section{Implementation}
\label{sec:impl}

\paragraph{Mask \fetch latencies by concurrent fetches and bounded staleness.}
The on-path retrieval of \fetch introduces latency penalties, ranging from microsecond to millisecond level when retrieving from geologically-distributed storage nodes.
Because the execution layer stalls until \fetch completes, the utilization is undermined, resulting in throughput degradation.
To recover the throughput, we introduce concurrent fetches.
The execution layer not only issue \fetch for the current transaction, but also pre-issues \fetch-es for the pending transactions in the execution queue.
By the time the transaction is ready to execute, its \fetch-ed values are likely already retrieved, allowing the transaction to proceed without stalling.

To perform concurrent fetches, the execution layer must predict the \fetch-es issued for the transactions before their proceeding transactions complete.
This is usually possible, and for certain common types of transactions (including UTXO ones), the \fetch-es can be completely determined by static analysis.
\gd{Remark that Ethereum smart contracts enforce explicit declaration of accessed keys.}
Nonetheless, for transactions whose \fetch-es cannot be (accurately) predicted, the execution layer can always issue on-path \fetch to execute the transactions.
\sys supports all kinds of transactions with an additional optimization for common types of transactions.

However, if the storage nodes can only respond to the retrieval requests for their current versions, they cannot respond to the concurrent \fetch-es before their trusted replica nodes \bump to the corresponding versions, and the sequential data dependency remains.
To address these, we introduce a bounded staleness semantic for \fetch.
The execution layer is configured with a staleness bound $s$, and it only issues \fetch-es for versions within $s$ of the current version.
At the same time, the storage nodes are configured with the same $s$, and they can respond to retrieval requests with values up to $s$ versions stale from the requested one.
Then, the execution layer further maintains a small in-memory cache of recent $s$ \bump operations.
When serving transaction reads, the cached updates supersede the \fetch-ed values if they overlap.
This cache not only guarantees to mask the staleness introduced by concurrent fetches but also potentially improves performance: if a hot key is frequently updated and fetched within a short period, the reads can be served directly from the cache and the performance is not affected by the \fetch at all.
\lijl{Is it beneficial to also cache a small set of non-hosting key-values for skewed workloads?} \gd{Mentioned.}

\paragraph{Designated pushing nodes.}
In \autoref{sec:design}, the \retrievev and \retrieveu requests are broadcast to all $r$ responsible storage nodes, and the first response is used.
This design amplifies the network traffic by $r$ times.
To reduce the network traffic, we can designate one storage node among the $r$ responsible nodes as the \emph{pushing node} for each version.
The pushing node actively broadcast \retrievev and \retrieveu responses to all replica nodes when they are available.
The replica nodes only send \retrievev and \retrieveu requests if they haven't received the corresponding responses from the pushing node after a timeout.

If the pushing node is faulty, this optimization may introduce additional latency due to the timeout.
To mitigate this, we order the $r$ responsible nodes deterministically, and a storage node only \emph{suppress} its pushing if it has observed that (one of) the preceding node pushes correctly for consecutive $K$ versions.
At the same time, the replica nodes send \retrievev and \retrieveu requests for $K$ versions after the pushing node times out.
This way, after the pushing node becomes faulty, the next correct responsible node takes over the pushing role immediately, and the faulty pushing node can only affect the latency of 1 version in every $K$ versions.
